% Section 2.1, from lectures/L02/appendix/a_io_ports.md.

\section{I/O Ports}
\label{c2:sec:ports}

\subsection{Three registers per port}
\label{c2:sec:three}
You have already used two of these. \code{sbi DDRB, 5} and \code{sbi PORTB, 5} lit an LED at the
end of Chapter~1 (\secref{c1:sec:program}), on one sentence's worth of explanation: \code{DDRx}
decides whether a pin is an output and \code{PORTx} decides what an output drives. That sentence
is true and it is about a third of the story, and the missing two thirds are what makes a driver
possible.

The ATmega328P has three I/O ports available on an Arduino Uno, called B, C and D, and each is
controlled by three registers. For port B those are \code{DDRB}, \code{PORTB} and \code{PINB}, at
data space addresses \code{0x24}, \code{0x25} and \code{0x23} (\secref{c1:sec:iowindow}).

The names are unhelpful in a specific way, so it is worth stating what each one is before anything
else:

\begin{booktable}{@{}p{4.4em}L>{\raggedright\arraybackslash}p{6.4em}@{}}
  \thead{Register} & \thead{What it is} & \thead{Read or write} \\
  \midrule
  \code{DDRx} & \textbf{D}ata \textbf{D}irection \textbf{R}egister. One bit per pin: 0 is input,
    1 is output. & Both \\
  \code{PORTx} & For an output pin, the level to drive. For an input pin, the pull-up switch.
    & Both \\
  \code{PINx} & The level the pin is actually at. Writing a 1 to it toggles \code{PORTx}.
    & Both, differently \\
\end{booktable}

\code{PORTx} is the one that trips people up, because it does two entirely unrelated jobs
depending on what \code{DDRx} says. And \code{PINx} is the one that surprises them later, because
writing to it does something that has nothing to do with reading it.

\subsection{Four states of a pin}
\label{c2:sec:states}
Two control bits, four combinations, and each one is a different piece of hardware behaviour.

\bookfigure{port_states}{\code{DDRx} against \code{PORTx}: a high-impedance input, an input with
its pull-up enabled, an output driven low, and an output driven high. \code{PINx} is the third
register and not part of the grid.}{c2:fig:states}

\textbf{A floating input is the dangerous one.} With \code{DDRx = 0} and \code{PORTx = 0} the pin
is connected to nothing at all, and reading \code{PINx} gives you whatever the surrounding
electrical noise induced on it. It is not random in a useful sense and it is not stable; it will
read 1 while your finger is near the board and 0 when it is not. An input you have not either
pulled up or wired to something is a bug, not a default.

\textbf{The pull-up is the reason the other three states are easy.} Setting \code{PORTx = 1} on an
input pin connects an internal resistor of 20 to 50\,k$\Omega$ between the pin and 5\,V. Now the
pin reads 1 unless something actively pulls it down, which is exactly what a button to ground
does.

\textbf{An output ignores the pull-up entirely.} Once \code{DDRx = 1}, \code{PORTx} means ``the
level to drive'', and the resistor is out of the picture.

\subsection{Writing to \code{PINx} toggles the output}
\label{c2:sec:toggle}
This is the one genuinely strange thing in this section, and it is worth being explicit about
because the name gives no hint of it.

Writing a \textbf{1} to a bit of \code{PINx} \textbf{toggles the corresponding bit of
\code{PORTx}}. Writing a 0 does nothing at all. Reading \code{PINx} is unaffected and still gives
the pin's actual level.

So storing a mask into \code{PINB} flips exactly the bits the mask selects, and leaves the rest
alone without having to read anything first. A toggle is one store, and the alternative is a load,
an exclusive-or and a store. It is faster and, more importantly, it cannot be interrupted halfway
through, which is a property Chapter~3 will make matter.

Two things this does \textbf{not} mean:
\begin{itemize}
  \item It does not toggle the pin directly. It toggles \code{PORTx}, so on an input pin it toggles
    the pull-up, which is occasionally useful and usually a mistake.
  \item It is not ``writing to the input register'', whatever the name suggests. Nothing you write
    to \code{PINx} is ever read back from \code{PINx}.
\end{itemize}

\subsection{An LED, and a resistor that is not optional}
\label{c2:sec:resistor}

\bookfigure{led_circuit}{A pin leaving the ATmega328P through a 220\,$\Omega$ series resistor into
an LED and then to ground. The chip boundary is drawn dashed.}{c2:fig:led}

Configure the pin as an output, drive it high, and current flows through the LED. Drive it low and
it stops. That is the whole circuit, and the driver in \secref{c2:sec:driver} is a way of saying
it in assembly.

The resistor sets the current. An LED is not a resistor: its voltage barely changes with current,
so without something to limit it the current is set by how much the pin can supply, which is about
40\,mA, which is roughly twice what a small LED survives and exactly the pin's absolute maximum.
The usual outcome is a dead LED; the occasional outcome is a dead pin.

An Arduino Uno already has an LED and a resistor on pin 13, which is why every example in this
book uses pin 13 and why none of them requires you to wire anything.

\textbf{The other way round works too, and you will meet it.} An LED can be wired from 5\,V
through its resistor to the pin, in which case driving the pin \emph{low} lights it. That is called
sinking rather than sourcing, and it inverts everything your driver means. It is not wrong; it is
a different board, and reading \code{led_on} without knowing which one you have is how you end up
with a program that is exactly backwards. On this part the two directions are specified the same,
whatever the folklore about older AVRs says: 20\,mA either way.

\subsection{A button, and why pressed reads zero}
\label{c2:sec:button}

\bookfigure{button_circuit}{A pin with the internal 20 to 50\,k$\Omega$ pull-up inside the chip
boundary, and a button to ground outside it. The pull-up is switched on by writing 1 to
\code{PORTB} with \code{DDRB} at 0.}{c2:fig:button}

Configure the pin as an input and switch the pull-up on. With the button not pressed, nothing else
is connected to the pin, and the pull-up holds it at 5\,V, so \code{PINx} reads \textbf{1}. Press
the button and the pin is shorted to ground, so \code{PINx} reads \textbf{0}.

\textbf{A pressed button reads zero.} Every driver in this book that reports ``is it pressed''
therefore has to invert what it read, and forgetting to is a bug that produces a program which does
exactly the opposite of what it should, reliably, which is at least easy to spot.

The alternative wiring, a button from the pin to 5\,V with an external pull-down resistor, gives
the intuitive polarity. Nobody does it, because it needs a resistor you have to fit and the
pull-up is free.

\textbf{What the pull-up does not do is debounce.} A mechanical button does not close once; it
closes, bounces open, closes again, several times over a few milliseconds. The pin faithfully
reports every one of those, so a program that counts presses by watching for a 1-to-0 transition
will count one press as four. That is not this chapter's problem, but it is the reason
Chapter~3's interrupt-driven button behaves in a way you would not predict from this circuit alone.

\subsection{Arduino pin numbers are not port bits}
\label{c2:sec:pinmap}

\bookfigure{arduino_pins}{Fourteen cells, one per Arduino digital pin: pins 0 to 7 are port D bits
0 to 7, and pins 8 to 13 are port B bits 0 to 5.}{c2:fig:pins}

The number silkscreened on the board is the board's idea. The device has never heard of it, and a
different board with the same chip numbers its pins differently.

The mapping is simple enough to state in one line, which is fortunate, because your driver has to
do it at run time:
\begin{itemize}
  \item Arduino pins \textbf{0 to 7} are port \textbf{D}, bits 0 to 7.
  \item Arduino pins \textbf{8 to 13} are port \textbf{B}, bits \textbf{0 to 5}.
\end{itemize}

The step at pin 8 is where implementations go wrong. Pin 8 is bit \textbf{0}, not bit 8, so the
driver has to subtract 8 before it shifts. Shifting by 8 in an eight-bit register gives zero
(\secref{c1:sec:shiftbits}), so the mask is zero, the LED never lights, nothing reports an error,
and the symptom looks like a wiring fault.

Port B stops at bit 5 because bits 6 and 7 are the crystal pins on an Arduino Uno. They exist on
the chip and are not available to you.

\textbf{Build this table yourself before reading it again.} Pin 13 is the one that matters here,
and deriving the other rows from the rule rather than reading them off is the difference between
knowing the mapping and having seen it.

\subsection{Reading, modifying, writing}
\label{c2:sec:rmw}
A port register holds eight pins. Your LED is one of them. So a driver that wants to turn its LED
on must not write to \code{PORTB}; it must change one bit of \code{PORTB} and leave the other seven
as it found them.

Three instructions and five cycles:
\begin{enumerate}
  \item \textbf{Read} the whole register into a scratch register, two cycles.
  \item \textbf{Modify} it, one cycle, by \code{or}-ing in the mask that selects your bit.
  \item \textbf{Write} the whole register back, two cycles.
\end{enumerate}

The middle step is where the mask earned its keep. To clear a bit instead, the mask is inverted and
the \code{or} becomes an \code{and}, which is why \code{shift_bits_inverted} exists as its own
routine rather than as a \code{com} after \code{shift_bits}.

A driver that assigns instead, \code{sts PORTB + 0x20, r24}, works perfectly with one LED and turns
off every other output on the port. It passes any test that only has one LED in it, which is why
the test suite has one with two.

\textbf{Read-modify-write is not atomic}, and on this machine that is three separate instructions
with two gaps in them. Nothing in this chapter can interrupt those gaps. From Chapter~3 onwards
something can, and that is where \code{sbi}, \code{cbi} and the \code{PINx} toggle stop being
micro-optimisations and start being correctness.
